% Preambolo condiviso dai documenti del progetto KeyItaly.
% Stile accademico sobrio: nessun colore decorativo, tipografia standard.

\usepackage[T1]{fontenc}
\usepackage[utf8]{inputenc}
\usepackage[italian]{babel}
\usepackage{lmodern}
\usepackage{microtype}

\usepackage[a4paper,margin=2.8cm]{geometry}
\usepackage{graphicx}
\usepackage{booktabs}
\usepackage{array}
\usepackage{enumitem}
\usepackage{xcolor}
\usepackage{tikz}
\usepackage{float}
\usepackage{pdflscape}
\usepackage[font=small,labelfont=bf,width=.92\textwidth]{caption}
\usepackage{listings}

\usepackage[hidelinks,pdfusetitle]{hyperref}
\usepackage{url}
\urlstyle{same}

\usetikzlibrary{positioning,shapes.geometric,arrows.meta,calc,fit,backgrounds}

% Interlinea leggermente aperta, più leggibile in stampa
\usepackage{setspace}
\onehalfspacing

% Elenchi più compatti dei default di LaTeX
\setlist[itemize]{itemsep=2pt,topsep=4pt,parsep=0pt}
\setlist[enumerate]{itemsep=2pt,topsep=4pt,parsep=0pt}
\setlist[description]{itemsep=4pt,topsep=6pt}

% ---------------------------------------------------------------------------
% Diagrammi
%
% La codifica cromatica è la stessa in tutto il documento:
%   grigio = accessibile a chiunque
%   azzurro = riservato all'utente registrato
%   rosso = riservato all'amministratore
%   giallo = vista JSP
%   verde = pattern di URL protetto da un filtro
% ---------------------------------------------------------------------------
\definecolor{cPubblico}{HTML}{ECECEC}
\definecolor{cUtente}{HTML}{CFE2F7}
\definecolor{cAdmin}{HTML}{F7D2D2}
\definecolor{cVista}{HTML}{FDEDC8}
\definecolor{cFiltro}{HTML}{D8EBD5}
\definecolor{cBordo}{HTML}{5A5A5A}

\tikzset{
  nodo/.style={
    draw=cBordo, line width=.5pt, rounded corners=2pt,
    align=center, font=\scriptsize, inner sep=2pt,
    minimum height=7mm, text width=20mm},
  pubblico/.style={nodo, fill=cPubblico},
  utente/.style={nodo, fill=cUtente},
  admin/.style={nodo, fill=cAdmin},
  vista/.style={nodo, fill=cVista},
  filtro/.style={nodo, fill=cFiltro},
  linea/.style={draw=cBordo, line width=.5pt},
  freccia/.style={linea, -{Stealth[length=4pt,width=3pt]}},
  tratteggio/.style={freccia, dashed, dash pattern=on 2.5pt off 2pt},
  % Notazione entità-relazione
  entita/.style={draw=black, line width=.6pt, fill=white,
    align=center, font=\small, minimum height=8mm, inner xsep=8pt},
  relazione/.style={draw=black, line width=.6pt, fill=white, diamond,
    aspect=2.2, align=center, font=\scriptsize, inner sep=1pt},
  composto/.style={draw=black, line width=.6pt, fill=white, ellipse,
    align=center, font=\scriptsize, inner sep=2pt},
  pallino/.style={draw=black, line width=.5pt, fill=white, circle,
    inner sep=0pt, minimum size=3.2pt},
  pallinoPieno/.style={pallino, fill=black},
  etichetta/.style={font=\scriptsize, inner sep=1.5pt},
}

% Attributo di un'entità nella notazione entità-relazione.
%   #1 stile del pallino (pallino oppure pallinoPieno)
%   #2 punto di attacco sul bordo dell'entità
%   #3 spostamento fino al pallino
%   #4 ancoraggio dell'etichetta rispetto al pallino
%   #5 testo dell'etichetta
\newcommand{\attributo}[5]{%
  \draw[line width=.5pt] (#2) -- ++(#3) node[#1] (ultimoAttr) {};
  \node[etichetta, anchor=#4, align=center] at (ultimoAttr) {#5};%
}

% Le figure occupano quasi tutta la larghezza del testo
\newcommand{\figuraLarga}[3]{%
  \begin{figure}[H]
    \centering
    \includegraphics[width=\linewidth]{#1}
    \caption{#2}\label{#3}
  \end{figure}}

% Stile per i frammenti di codice e configurazione
\lstdefinestyle{sobrio}{
  basicstyle=\ttfamily\footnotesize,
  keywordstyle=\bfseries,
  commentstyle=\itshape,
  numbers=none,
  frame=single,
  rulecolor=\color{black!35},
  framesep=6pt,
  breaklines=true,
  showstringspaces=false,
  captionpos=b,
  extendedchars=true,
  literate={à}{{\`a}}1 {è}{{\`e}}1 {é}{{\'e}}1 {ì}{{\`i}}1 {ò}{{\`o}}1 {ù}{{\`u}}1
}
\lstset{style=sobrio}

% Frontespizio
\newcommand{\frontespizio}[2]{%
  \begin{titlepage}
    \centering
    \vspace*{2cm}
    {\LARGE Università degli Studi di Salerno\par}
    \vspace{4mm}
    {\large Dipartimento di Informatica\par}
    {\large Corso di Laurea in Informatica\par}
    \vspace{2.2cm}
    {\large Tecnologie Software per il Web\par}
    \vspace{1.4cm}
    \rule{\linewidth}{0.4pt}\par
    \vspace{6mm}
    {\huge\bfseries #1\par}
    \vspace{4mm}
    {\Large\itshape KeyItaly\par}
    \vspace{3mm}
    {\large #2\par}
    \vspace{6mm}
    \rule{\linewidth}{0.4pt}\par
    \vfill
    \begin{minipage}[t]{0.45\textwidth}
      \raggedright
      \textbf{Docente}\par
      \vspace{2mm}
      Prof.ssa Rita Francese
    \end{minipage}%
    \begin{minipage}[t]{0.45\textwidth}
      \raggedleft
      \textbf{Studenti}\par
      \vspace{2mm}
      Giovanni Cerchia\quad 0512120818\par
      Andrea Compagnone\quad 0512116159
    \end{minipage}
    \vspace{1.4cm}
    \par
    {\large Anno Accademico 2023/2024\par}
    \vspace{1cm}
  \end{titlepage}}
